% GENERATED FILE. Edit canonical lessons, then run npm run generate:books.
% canonical-source-hash: fnv1a64:260a0b3f6aa74bd8
% canonical-lessons: ZH-C121-jiao-teach, ZH-C121-lianxi, ZH-C121-huida, ZH-C121-fuxi, ZH-C121-cuo

\chapter{Teaching and Practising}
\label{ch:zh-teaching-and-practising}
\hlchaptermodality{121}

\begin{chapteropening}
\textbf{By the end of this chapter:} \emph{I can say teach, practise, answer and revise, and say something is wrong.}

Complete the last lesson of chapter 121: I can say teach, practise, answer and revise, and say something is wrong.
\end{chapteropening}

\section[jiāo]{jiāo — to teach}

\label{lesson:ZH-C121-jiao-teach}



\begin{quote}
Before the new one: say the Mandarin for to pack food to take away, then the Mandarin for dim sum.
\end{quote}

\subsection*{You'll want to know: jiāo}
\emph{jiāo} — \textquotedblleft{}to teach\textquotedblright{}.

\begin{quote}
\emph{tā jiāo wǒ hànyǔ} — She teaches me Chinese.
\end{quote}

Its characters come later; for now it is a word for the ear.

\begin{cousinweb}
\emph{jiāo} is inherited from Middle Chinese \emph{kaewH} (Guangyun, in Baxter's transcription), glossed 'to teach'; the dictionary also lists the character with a level tone, \emph{kaew}, glossed 'to imitate', and the level jiāo of the verb matches that one.
\end{cousinweb}

\begin{grammarlens}[title={What you've built}]
One more word for everyday actions.
\end{grammarlens}

\subsection*{Your turn}
\begin{itemize}
\raggedright
  \item \emph{Say it:} \emph{jiāo}
  \item \emph{Say it:} \emph{jiāo}, once more
  \item \emph{Say it:} \emph{tā jiāo wǒ hànyǔ} — She teaches me Chinese.
  \item \emph{Recall:} say the Mandarin for to pack food to take away, then the Mandarin for dim sum, then say \emph{jiāo} again
\end{itemize}

\begin{tcolorbox}[breakable,colback=teal!4,colframe=teal!35!black,title={Before you move on}]
What does jiāo mean? (\textquotedblleft{}To teach\textquotedblright{}.) Say it once more.
\end{tcolorbox}
\section[liànxí]{liànxí — to practise}

\label{lesson:ZH-C121-lianxi}



\begin{quote}
Before the new one: say the Mandarin for dim sum, then the Mandarin for to teach.
\end{quote}

\subsection*{You'll want to know: liànxí}
\emph{liànxí} — \textquotedblleft{}to practise\textquotedblright{}.

\begin{quote}
\emph{wǒ měitiān liànxí shuō hànyǔ} — I practise speaking Chinese every day.
\end{quote}

Its characters come later; for now it is a word for the ear.

\begin{cousinweb}
\emph{liànxí} is built from liàn and xí, both inherited from Middle Chinese, \emph{lenH} and \emph{zip} (Guangyun, in Baxter's transcription); the dictionary glosses the first as white boiled silk and the second 'to learn', explaining it as flying again and again.
\end{cousinweb}

\begin{grammarlens}[title={What you've built}]
One more word for everyday actions.
\end{grammarlens}

\subsection*{Your turn}
\begin{itemize}
\raggedright
  \item \emph{Say it:} \emph{liànxí}
  \item \emph{Say it:} \emph{liànxí}, once more
  \item \emph{Say it:} \emph{wǒ měitiān liànxí shuō hànyǔ} — I practise speaking Chinese every day.
  \item \emph{Recall:} say the Mandarin for dim sum, then the Mandarin for to teach, then say \emph{liànxí} again
\end{itemize}

\begin{tcolorbox}[breakable,colback=teal!4,colframe=teal!35!black,title={Before you move on}]
What does liànxí mean? (\textquotedblleft{}To practise\textquotedblright{}.) Say it once more.
\end{tcolorbox}
\section[huídá]{huídá — to answer}

\label{lesson:ZH-C121-huida}



\begin{quote}
Before the new one: say the Mandarin for to teach, then the Mandarin for to practise.
\end{quote}

\subsection*{You'll want to know: huídá}
\emph{huídá} — \textquotedblleft{}to answer\textquotedblright{}.

\begin{quote}
\emph{qǐng nǐ huídá} — Please answer.
\end{quote}

Its characters come later; for now it is a word for the ear.

\begin{cousinweb}
\emph{huídá} is built from huí 'to turn back' and dá 'to respond', both inherited from Middle Chinese, \emph{hwoj} and \emph{top} (Guangyun, in Baxter's transcription); the dictionary glosses the second 'to match'.
\end{cousinweb}

\begin{grammarlens}[title={What you've built}]
One more word for everyday actions.
\end{grammarlens}

\subsection*{Your turn}
\begin{itemize}
\raggedright
  \item \emph{Say it:} \emph{huídá}
  \item \emph{Say it:} \emph{huídá}, once more
  \item \emph{Say it:} \emph{qǐng nǐ huídá} — Please answer.
  \item \emph{Recall:} say the Mandarin for to teach, then the Mandarin for to practise, then say \emph{huídá} again
\end{itemize}

\begin{tcolorbox}[breakable,colback=teal!4,colframe=teal!35!black,title={Before you move on}]
What does huídá mean? (\textquotedblleft{}To answer\textquotedblright{}.) Say it once more.
\end{tcolorbox}
\section[fùxí]{fùxí — to revise; to go over again}

\label{lesson:ZH-C121-fuxi}



\begin{quote}
Before the new one: say the Mandarin for to practise, then the Mandarin for to answer.
\end{quote}

\subsection*{You'll want to know: fùxí}
\emph{fùxí} — \textquotedblleft{}to revise; to go over again\textquotedblright{}.

\begin{quote}
\emph{wǒ wǎnshang fùxí hànyǔ} — I go over my Chinese in the evening.
\end{quote}

Its characters come later; for now it is a word for the ear.

\begin{cousinweb}
\emph{fùxí} is built from fù 'again' and xí 'to learn', both inherited from Middle Chinese, \emph{bjuwk} and \emph{zip} (Guangyun, in Baxter's transcription): to learn again. The dictionary glosses the first 'to return, to repeat' and lists it with a departing tone too, \emph{bjuwH}, glossed 'again': fù continues that one, since the regular change from \emph{bjuwk} gives fú.
\end{cousinweb}

\begin{grammarlens}[title={What you've built}]
One more word for everyday actions.
\end{grammarlens}

\subsection*{Your turn}
\begin{itemize}
\raggedright
  \item \emph{Say it:} \emph{fùxí}
  \item \emph{Say it:} \emph{fùxí}, once more
  \item \emph{Say it:} \emph{wǒ wǎnshang fùxí hànyǔ} — I go over my Chinese in the evening.
  \item \emph{Recall:} say the Mandarin for to practise, then the Mandarin for to answer, then say \emph{fùxí} again
\end{itemize}

\begin{tcolorbox}[breakable,colback=teal!4,colframe=teal!35!black,title={Before you move on}]
What does fùxí mean? (\textquotedblleft{}To revise; to go over again\textquotedblright{}.) Say it once more.
\end{tcolorbox}
\section[cuò]{cuò — wrong; mistaken}

\label{lesson:ZH-C121-cuo}



\begin{quote}
Before the new one: say the Mandarin for to answer, then the Mandarin for to revise.
\end{quote}

\subsection*{You'll want to know: cuò}
\emph{cuò} — \textquotedblleft{}wrong; mistaken\textquotedblright{}.

\begin{quote}
\emph{zhège hànzì cuò le} — This character is wrong.
\end{quote}

Its characters come later; for now it is a word for the ear.

\begin{cousinweb}
\emph{cuò} is inherited from Middle Chinese \emph{tshak} (Guangyun, in Baxter's transcription), glossed 'mixed' and 'crossing slantwise'; how it came to mean wrong is not settled.
\end{cousinweb}

\begin{grammarlens}[title={What you've built}]
One more word for everyday actions.
\end{grammarlens}

\subsection*{Your turn}
\begin{itemize}
\raggedright
  \item \emph{Say it:} \emph{cuò}
  \item \emph{Say it:} \emph{cuò}, once more
  \item \emph{Say it:} \emph{zhège hànzì cuò le} — This character is wrong.
  \item \emph{Recall:} say the Mandarin for to answer, then the Mandarin for to revise, then say \emph{cuò} again
\end{itemize}

\begin{tcolorbox}[breakable,colback=teal!4,colframe=teal!35!black,title={Before you move on}]
What does cuò mean? (\textquotedblleft{}Wrong; mistaken\textquotedblright{}.) Say it once more.
\end{tcolorbox}
